\documentclass[10pt,twocolumn,letterpaper]{article}
\usepackage{balance} % in the preamble
\usepackage{cvpr}
\usepackage{times}
\usepackage{epsfig}
\usepackage{graphicx}
\usepackage{amsmath}
\usepackage{amssymb}
\usepackage{comment}

% Include other packages here, before hyperref.
\usepackage[breaklinks=true,bookmarks=false]{hyperref}

\cvprfinalcopy % *** Uncomment this line for the final submission

\def\cvprPaperID{****} % *** Enter the CVPR Paper ID here
\def\httilde{\mbox{\tt\raisebox{-.5ex}{\symbol{126}}}}

\begin{document}

%%%%%%%%% TITLE
\title{PROJECT PROPOSAL: Comparing Dimensionality Reduction and Classifier Architectures for MultiROCKET-Based ECG Classification: A Controlled Ablation Study}

\author{Ramita Rathore\\
Georgia State University\\
Atlanta, GA, USA \\
{\tt\small rrathore1@student.gsu.edu}
\and
Jacob Cartwright\\
Georgia State University\\
Atlanta, GA, USA \\
{\tt\small jcartwright12@student.gsu.edu}
\and
Hussain Manasi\\
Georgia State University\\
Atlanta, GA, USA \\
{\tt\small hmanasi1@student.gsu.edu}
}

\maketitle

%%%%%%%%% BODY TEXT
\section{Introduction}

Automated electrocardiogram (ECG) interpretation can assist clinical decision-making by providing rapid, consistent diagnoses, particularly in settings with limited access to specialists. However, multi-label ECG classification remains challenging due to the co-occurrence of multiple cardiac conditions, severe class imbalance between common and rare diagnoses, and the need to capture complex spatiotemporal patterns across twelve leads.

Recent work has shown that MultiROCKET—a random convolutional feature transform—combined with simple linear classifiers achieves competitive performance at a fraction of the computational cost of deep learning \cite{tan2022, wang2023}. However, the raw MultiROCKET feature space is extremely high-dimensional (often exceeding 16,000 features), noisy, and contains many redundant components. How best to reduce this feature space, and which downstream classifier architecture is most effective, remains an open question.

This paper investigates two orthogonal design choices in a controlled $2 \times 2$ ablation: (1) \textbf{dimensionality reduction strategy}—either a nonlinear autoencoder or structured kernel/statistic/origin pooling—and (2) \textbf{classifier architecture}—either a shallow MLP or an FT-Transformer. We evaluate all four combinations on a large 12-lead ECG dataset, focusing on a four-class diagnostic label subset selected to enable direct comparison against prior benchmark results, as reported by \cite{zheng2022}.

The central contributions of this work are: (i) a rigorous comparison of compression strategies for MultiROCKET features; (ii) an empirical evaluation of whether transformer-based tabular models (FT-Transformer) outperform MLPs on dense real-valued feature inputs; and (iii) reproducible benchmarking against published results on a shared label subset.


\section{Related Works}

The evolution of automated 12-lead ECG interpretation has shifted from traditional signal processing toward integrated frameworks that combine high-fidelity data repositories, efficient feature extraction, and modular deep learning architectures.

\subsection{Large-Scale 12-Lead ECG Data}

The transition from small, single-lead datasets to massive multivariate repositories has been critical for model generalization. Zheng et al.\ introduced a landmark 12-lead database comprising 45,152 patients sampled at 500 Hz \cite{zheng2022}. Unlike the legacy MIT-BIH database, which is limited by its small patient cohort \cite{frontiers2022}, the Zheng et al.\ repository adheres to AHA/ACC/HRS recommendations, mapping 56 cardiovascular conditions to standardized SNOMED-CT codes \cite{kurias2023}. This granularity allows for robust multi-label classification across diverse demographics \cite{zheng_multilabel}.

\subsection{The ROCKET Family for Feature Extraction}

While conventional deep learning models like CNNs and LSTMs excel at spatial patterns, they are often computationally intensive over long ECG horizons \cite{kth_minirocket, nuyen2024}. Dempster et al.\ proposed ROCKET as a paradigm shift, utilizing random convolutional kernels to achieve state-of-the-art accuracy with minimal training time \cite{dempster2020}. A key innovation is the ``Proportion of Positive Values'' (PPV) feature, which effectively captures chaotic irregularity characteristic of arrhythmias such as atrial fibrillation \cite{arxiv_calf}.

Subsequent iterations have optimized this approach: \textbf{MiniRocket} provides a deterministic transform for millisecond-level inference on edge devices \cite{dempster2021}, while \textbf{MultiRocket} incorporates first-order derivatives and multiple pooling operators to capture nuanced wave morphology and slope dynamics \cite{tan2022, wang2023}.

\subsection{Dimensionality Reduction for High-Dimensional Feature Spaces}

A significant challenge in 12-lead analysis is the ``Curse of Dimensionality,'' as lead-wise ROCKET transforms can produce over 600,000 features \cite{researchgate_block}. Uribarri et al.\ addressed this with \textbf{Detach-ROCKET}, which utilizes Sequential Feature Detachment (SFD) to iteratively prune low-contribution features \cite{uribarri2023}, compressing the feature space by up to 98\% without degrading diagnostic accuracy \cite{upcommons_rocket}.

An alternative reduction strategy is the autoencoder, which performs nonlinear compression of high-dimensional inputs into a compact latent representation. Variations of autoencoders have been applied in ECG contexts to denoise signals \cite{nurmaini_deep_2020, xia_denoising_2023} and learn compact representations prior to downstream classification \cite{c_harvey_ecg_2025, g_wang_unsupervised_2023}. In contrast to feature pruning methods like Detach-ROCKET, the encoder produces a learned dense embedding that may preserve richer nonlinear structure—though at the cost of an additional training stage.

A simpler alternative is \textbf{structured pooling}. In our implementation, we use \textbf{dilation-aware} pooling: the fitted MultiROCKET transformer provides per-kernel metadata (dilation, statistic, origin). Features are grouped by (dilation, statistic, origin); within each group, mean pooling is applied, yielding one value per group (e.g., 22 dilations $\times$ 4 stats $\times$ 2 origins $=$ 176 dimensions). This approach requires no learned parameters and preserves interpretable structure (each dimension is a specific dilation--statistic--origin triple).

\subsection{Transformer Models for Tabular Data}

The FT-Transformer \cite{gorishniy_revisiting_2023} adapts the transformer architecture to tabular (structured, real-valued) data by treating each input feature as a token and applying self-attention across the feature dimension. This allows the model to explicitly capture pairwise feature interactions—for example, ``feature A high AND feature B low implies class X''—which may be clinically relevant in ECG diagnosis where conditions often depend on combinations of morphological and rhythm-based findings. Gorishniy et al.\ demonstrated that FT-Transformer is competitive with or superior to gradient boosting methods on a range of tabular benchmarks, though comparisons with MLPs on purely real-valued (non-categorical) inputs are less conclusive \cite{gorishniy_revisiting_2023}.

\subsection{Benchmarking ECG Classification}

To situate our results within the broader literature, we adopt the four diagnostic label subset used by \cite{zheng2022}, which includes atrial fibrillation, general supraventricular tachycardia, sinus bradycardia and sinus rhythm. Reporting on this shared subset enables direct numerical comparison against confusion matrices and weighted average F-1 scores, providing an external validity check beyond internal ablation results \cite{zheng2022}.


\section{Planned Methodology}

\subsection{Dataset and Problem Setting}

We use the expanded Chapman-Shaoxing dataset \cite{zheng2022}. After loading and validation, 45,150 single 10-second ECG recordings remain, each sampled at 500 Hz across 12 leads (5,000 samples per lead); shape $(45{,}150 \times 12 \times 5{,}000)$. Following \cite{zheng2022}, we restrict the target label space to four diagnostic conditions: atrial fibrillation (AF), general supraventricular tachycardia (SVT), sinus bradycardia, and sinus rhythm. Labels are encoded as a single class per recording (multiclass). This focused label subset enables direct benchmarking against published results while controlling label imbalance complexity.

The dataset is partitioned into train/validation/test sets using a 60/20/20 ratio with stratified random sampling (seed 0). Resulting split sizes: approximately 27,090 train, 9,030 validation, 9,030 test. Split indices are generated once and serialized to disk (\texttt{splits.npz}), ensuring identical partitions across all experiments and machines.

\subsection{Core Design: A $2 \times 2$ Ablation}

Our study is structured as a fully crossed $2 \times 2$ ablation over two independent design dimensions:

\begin{itemize}
    \item \textbf{Dimensionality Reduction:} Autoencoder (Path A) vs.\ Structured Pooling (Path B)
    \item \textbf{Classifier:} MLP vs.\ FT-Transformer
\end{itemize}

This yields four experimental conditions, summarized in Table~\ref{tab:conditions}. All four conditions share an identical upstream pipeline: raw 12-lead ECG $\rightarrow$ MultiROCKET feature extraction $\rightarrow$ feature scaling. Downstream performance differences are therefore attributable solely to the reduction strategy and classifier architecture.

\begin{table}[h]
\centering
\small
\begin{tabular}{lcc}
\hline
\textbf{Condition} & \textbf{Reduction} & \textbf{Classifier} \\
\hline
A1 & Autoencoder & MLP \\
A2 & Autoencoder & FT-Transformer \\
B1 & Structured Pooling & MLP \\
B2 & Structured Pooling & FT-Transformer \\
\hline
\end{tabular}
\caption{The four experimental conditions in the $2 \times 2$ ablation.}
\label{tab:conditions}
\end{table}

\subsection{Feature Extraction}

All experiments use a fixed feature representation computed with \emph{multivariate MultiROCKET} (sktime \texttt{MultiRocketMultivariate}), fitted on the training set only. The transformer uses 2,048 random kernels and outputs $2 \times 4 \times 2{,}016 = 16{,}128$ features per recording: 2 signal origins (raw and first-order difference), 4 statistics per origin (PPV, LSPV, MPV, MIPV\footnote{Proportion of Positive Values; Longest Stretch of Positive Values; Max; Mean of Positive Values.}), and 2,016 base feature dimensions determined by the kernel layout (84 kernels per group $\times$ 24 feature multipliers across dilations).

Features are computed in batches of 1,024 and stored as disk-backed \texttt{memmap} arrays in \texttt{float32} (\texttt{train.dat}, \texttt{test.dat}, \texttt{val.dat}) with shapes recorded in \texttt{shapes.npz}. The same fixed features are reused across all four experimental conditions without recomputation.

\subsection{Dimensionality Reduction}

\paragraph{Path A — Autoencoder.}
Raw MultiROCKET features (16,128 dimensions) are standardized with \texttt{StandardScaler} fit on the training set, then transformed on train/validation/test. A shallow feedforward autoencoder is trained to reconstruct these scaled vectors: input dimension 16,128, encoder hidden dimensions [512, 256] with ReLU, latent dimension $d = 256$. The decoder is symmetric (256 $\to$ 256 $\to$ 512 $\to$ 16,128). Training is unsupervised (MSE reconstruction loss), Adam optimizer (learning rate $10^{-3}$), batch size 256, for 50 epochs. Only the encoder is retained for downstream use; the classifier receives 256-dimensional latent vectors. The autoencoder is trained independently of the downstream classifier.

\paragraph{Path B — Structured Pooling.}
Kernel metadata from the fitted MultiROCKET transformer (dilations, \texttt{num\_features\_per\_dilation}, 84 kernels per group) is used to build dilation-aware pool groups. Features are grouped by \textbf{dilation} (22 levels), \textbf{statistic} (PPV, LSPV, MPV, MIPV), and \textbf{origin} (raw, differenced). Within each (dilation, statistic, origin) group, \textbf{mean} pooling is applied over the corresponding column indices, yielding one value per group: $22 \times 4 \times 2 = 176$ features. The pipeline is: (1) pre-pool \texttt{StandardScaler} on the 16,128 raw features (fit on train); (2) mean pooling over the 176 groups; (3) post-pool \texttt{StandardScaler} on the 176 pooled features (fit on train). No learned parameters; fully deterministic; each output dimension is a specific (dilation, statistic, origin) triple.

\subsection{Classifier Architectures}

\paragraph{MLP.}
A multi-layer perceptron: linear $\to$ ReLU $\to$ dropout (0.1) $\to$ linear $\to$ ReLU $\to$ dropout (0.1) $\to$ linear to 4 logits. Hidden dimensions [512, 256]. Output is a four-way probability distribution via softmax; trained with cross-entropy loss. Optimizer: AdamW, learning rate $10^{-3}$, weight decay $10^{-4}$, batch size 256, 20 epochs.

\paragraph{FT-Transformer.}
The Feature Tokenizer + Transformer architecture \cite{gorishniy_revisiting_2023} treats each of the $n$ reduced features as a token: each scalar is projected by a shared \texttt{Linear(1, d\_token)} (here $d_{\mathrm{token}}=64$), a [CLS] token is prepended, and a transformer encoder (2 layers, 4 heads, GELU, dropout 0.1) is applied. The classification head is a linear layer from the [CLS] representation to 4 logits (four-way multiclass, cross-entropy loss). Optimizer: AdamW, learning rate $2 \times 10^{-3}$, weight decay $10^{-4}$, batch size 1,024, 20 epochs. This architecture can explicitly model feature interactions. The instructor-identified concern, that FT-Transformer may not offer advantages over MLP on purely real-valued inputs, is addressed empirically by conditions A1 vs.\ A2 and B1 vs.\ B2.

\subsection{Evaluation}

All four conditions are evaluated on the same held-out test set under identical label definitions matching \cite{zheng2022}.

\paragraph{Primary metrics:}

Following the evaluation approach used by Zheng et al.\ \cite{zheng2022}, we adopt the following metrics:
\begin{itemize}
    \item Confusion and Normalized Confusion Matrices to evaluate models' classification performance
    \item Weighted Average F1-Score for model comparison and
    hyperparameter selection.
\end{itemize}

\paragraph{Ablation analysis:}
\begin{itemize}
    \item \textbf{Reduction effect:} Compare (A1+A2) vs.\ (B1+B2) to isolate the contribution of the autoencoder vs.\ structured pooling, averaged across classifiers.
    \item \textbf{Classifier effect:} Compare (A1+B1) vs.\ (A2+B2) to isolate MLP vs.\ FT-Transformer performance, averaged across reduction strategies.
    \item \textbf{Interaction effect:} Assess whether one classifier benefits disproportionately from one reduction strategy.
\end{itemize}

\paragraph{Benchmarking:}
Results on the four-class label subset are directly compared against the published figures of \cite{zheng2022} to contextualize performance relative to the broader literature.


\subsection{Summary}

This methodology holds the upstream feature representation constant and varies exactly two binary factors—reduction strategy and classifier architecture—enabling clean causal attribution of performance differences. Evaluating on a shared four-class label subset ensures external comparability. The $2 \times 2$ design answers three distinct research questions with only four experimental conditions, making it computationally tractable while yielding interpretable conclusions.


\section{Preliminary Report Goals}

The preliminary report will establish the complete data pipeline and baseline performance benchmarks. Specifically, we plan on delivering:

\textbf{Complete Sections:} Finalized Introduction and Related Work sections formatted per IEEE CVPR guidelines, including properly cited references.

\textbf{Dataset Preparation:} Preprocessed 12-lead ECG dataset with documented train/val/test splits. Exploratory analysis showing label distributions, class imbalance statistics, and co-occurrence patterns for the four target diagnostic labels as defined by \cite{zheng2022}.

\textbf{Feature Extraction:} Extracted multivariate MultiROCKET features stored as disk-backed arrays for efficient reuse across all four experimental conditions.

\textbf{Baseline Results:} At minimum, conditions B1 (Structured Pooling + MLP) and A1 (Autoencoder + MLP) trained and evaluated, with confusion matrices and weighted average F-1 scores reported on the four-class label subset. We will also compare our results to the results presented in the paper \cite{zheng2022} using regular and normalized confusion matrices to evaluate classification models, and weighted average F1-Score to select hyper parameters. 



\textbf{Figures \& Tables:} Sample ECG visualizations, class distribution plots for the four target labels, condition comparison tables, and performance breakdowns—all with proper captions and formatting.

The final report will add FT-Transformer conditions (A2, B2), complete the $2 \times 2$ ablation analysis, and provide full benchmarking against prior work.

%-------------------------------------------------------------------------
% REFERENCES SECTION
{
\small
\raggedright
\bibliographystyle{ieee_fullname}
\bibliography{egbib}
}

\end{document}